% ISSUE 2 — Supp. Table 4 (30/30) vs. Abstract/Conclusion (27/30) on HyperSymLoop (M4)
% Category: Stale (classic same-run desync)
% Source A: supp_benchmark_report.tex, Abstract, l.128-156
% Source B: supp_benchmark_report.tex, tab:overall (Six-Method Aggregate Performance), l.390-423

%% ---- SOURCE A: Abstract ----
\begin{abstract}
We present an extended evaluation of the HypatiaX symbolic regression
framework on 30 Feynman benchmark equations, integrating a six-method
noiseless comparison with a comprehensive noise robustness sweep
($\sigma \in \{0, 0.5, 1, 5, 10\}\%$) and a sample complexity sweep
($n \in \{50, 100, 200, 500, 750, 1000\}$), focusing on the two
top-performing core methods: \EHD{} (M3) and \HSL{} (M4).

Following correction of a measurement harness bug in \HSL{} (see
Section~\ref{sec:bugfix}), both methods achieve median $\Rsq = 1.000000$
across all tested conditions.
\EHD{} achieves 100\% recovery at all noise levels tested.
\HSL{} achieves 100\% recovery at all \emph{noisy} conditions ($\sigma > 0$)
and 90.0\% (27/30 equations) under the strict noiseless protocol ($\Rsq > 0.9999$).
All results in this document use the v2 corrected harness; v1 figures are
retained in the repository commit history for audit purposes only.
Both converge at $n = 50$ samples.
\EHD{} offers exact symbolic interpretability; \HSL{} is $1.5$--$76\times$
faster depending on noise level.

On this 30-equation subset, \EHD{} achieves 96.7\% recovery vs.\ AI~Feynman~2.0's
79.3\% ($+17.4$\,pp), and \HDS{} achieves 90.0\% ($+10.7$\,pp).
These are subset-specific comparisons: published baselines were evaluated on
all 100 Feynman equations under the standardised SRBench protocol; these
figures should not be read as a general state-of-the-art claim.
We additionally correct four methodological issues:
undisclosed hardcoded \PureLLM{} results, scale-biased raw \rmse,
missing Wilcoxon significance tests, and mismatched pass-rate thresholds.
\end{abstract}

%% ---- SOURCE B: tab:overall ----
%% ============================================================
\section{Six-Method Noiseless Results}
\label{sec:noiseless}

\begin{table}[ht]
\caption{Six-method aggregate performance, noiseless protocol
  ($\Rsq \geq 0.9999$, $n=1000$).
  \warn~\PureLLM{} pass rate includes 11/30 hardcoded results.}
\label{tab:overall}
\centering
\small
\renewcommand{\arraystretch}{1.3}
\begin{tabular}{L{2.4cm} C{1.0cm}
                C{1.8cm} C{1.3cm} C{1.4cm}}
\toprule
\textbf{Method} & \textbf{Type} &
  \makecell{\textbf{Pass Rate}\\$(\Rsq \geq 0.9999)$} &
  \makecell{\textbf{Mean} $\Rsq$} &
  \makecell{\textbf{Avg}\\Runtime} \\
\midrule
\rowcolor{warnAmberBg}
\PureLLM\,\warn & Core  & 30/30 (100\%)\warn & 1.000000 &   5.7\,s \\
\rowcolor{passgreenBg}
\EHD\,(M3)      & Core  & 29/30 (96.7\%)      & 0.999993 &  20.2\,s \\
\HDS            & Tools & 27/30 (90.0\%)      & 0.999960 & 155.0\,s \\
\HSL\,(M4)      & Core  & 30/30 (v2)$^\dagger$ & 0.999+  &  11.4\,s \\
\SEL            & Tools & 26/30 (86.7\%)      & 0.999829 & 201.2\,s \\
\rowcolor{failredBg}
\INN            & Core  & 17/30 (56.7\%)      & 0.999272 &  23.5\,s \\
\bottomrule
\multicolumn{5}{l}{\footnotesize $^\dagger$\,\HSL{} v1 showed 26/30 due to the Newton
  measurement bug; v2 achieves 30/30 (Section~\ref{sec:bugfix}).}\\
\end{tabular}
\end{table}
